\documentclass[11pt]{article}

\usepackage[margin=1in]{geometry}
\usepackage{amsmath,amsthm,amssymb}
\usepackage{mathtools}
\usepackage{booktabs}
\usepackage{array}
\usepackage{microtype}
\usepackage{hyperref}
\usepackage[capitalize,noabbrev]{cleveref}
\usepackage[
  backend=biber,
  style=numeric-comp,
  sorting=none
]{biblatex}
\addbibresource{refs.bib}

\newtheorem{theorem}{Theorem}
\newtheorem{lemma}[theorem]{Lemma}
\newtheorem{corollary}[theorem]{Corollary}
\theoremstyle{definition}
\newtheorem{definition}[theorem]{Definition}
\theoremstyle{remark}
\newtheorem*{remark}{Remark}

\title{Zero-Copy Wire Formats and the End of the Serialization Tax:\\
A Latency Model and Benchmark Methodology for the ZAP Protocol}
\author{ZAP Protocol Authors\\\texttt{zap-proto.io}}
\date{2026}

\begin{document}
\maketitle

\begin{abstract}
Modern microservice deployments spend a measurable fraction of CPU
on encoding and decoding RPC payloads. Profiling of warehouse-scale
workloads attributes 4--7\% of cycles to protocol-buffer-style
serialization in hot paths~\cite{kanev2015profiling}.
We present a formal latency model and a benchmark methodology for the
ZAP wire format, a Cap'n~Proto--derived encoding that places no
parsing step between received bytes and program-visible objects.
We decompose RPC latency into serialization, transit, deserialization,
and handler components, and show that for ZAP the serialization and
deserialization terms collapse to constant time in field count, leaving
only payload-byte cost. We give a throughput-bound theorem proving that
ZAP's effective bandwidth approaches link rate under bounded payload
size, while protobuf-over-HTTP/2 plateaus at roughly 60--70\%\ due to
HPACK and tag-length-value overhead. We describe a reproducible
single-machine and 10\,GbE benchmark methodology that reports
$p50/p95/p99/p999$ tails under both quiescent and congested regimes.
We discuss honestly when ZAP does \emph{not} win, including
single-field tiny messages and embedded environments without
zero-copy mmap.
\end{abstract}

\section{Introduction}
\label{sec:intro}

The serialization tax is the cost a service pays, measured in CPU
cycles and tail latency, to translate an in-memory object graph into a
sequence of bytes for transit and back. Profiling of large fleets shows
this is not a rounding error: in a warehouse-scale measurement,
serialization libraries together account for 4--7\%\ of all
cycles~\cite{kanev2015profiling}. At the 99.9th percentile, where users
notice, the cost is more pronounced because allocator pressure,
copy-induced cache misses, and HPACK dynamic-table churn each
contribute their own tail.

We argue, and prove, that this tax is not fundamental. It is an
artifact of designing wire formats around streaming parsers when the
deployment reality is mmap'd memory, persistent connections, and
schemas known at compile time. The ZAP wire format is derived from
Cap'n Proto~\cite{capnproto,capnproto-encoding}: messages are arenas
of fixed-layout structs connected by relative pointers, and parsing
collapses to the identity function on the byte buffer.

This paper makes three contributions:
\begin{enumerate}
  \item A formal RAM-model latency decomposition that separates
        per-field cost from per-byte cost.
  \item A throughput bound (\Cref{thm:throughput}) and a tail-latency
        independence result (\Cref{thm:tail}) for the ZAP wire format.
  \item A benchmark methodology covering small/medium/large RPC sizes,
        loopback and 10\,GbE, post-quantum handshake amortization, and
        honest discussion of regimes where ZAP loses.
\end{enumerate}

\section{Wire Format Recap}
\label{sec:wire}

\subsection{Arenas and segments}
A ZAP message is a sequence of \emph{segments}, each a contiguous byte
range obtained from an arena allocator. Segments are 8-byte aligned.
Multiple segments form a single logical message and may be transmitted
back-to-back, mmap'd from a file, or constructed in shared memory.

\subsection{Far-pointer model}
Pointers are 64-bit relative offsets, not virtual addresses. A pointer
encodes target segment id and word offset within the segment, plus a
small tag distinguishing struct, list, far, and capability pointers.
Because offsets are relative, no fixup is required when the message is
copied, mmap'd at a different address, or cast directly out of a
network buffer.

\subsection{Capability tables vs schema-on-the-wire}
Unlike formats that embed field tags in every record (Protocol
Buffers~\cite{protobuf}, MessagePack~\cite{msgpack}), ZAP places
schema information in compile-time-generated accessors. The on-wire
record is purely positional: byte $k$ of a struct of type $T$ is
field $T.f_k$ by construction. This is the same trade-off as
FlatBuffers~\cite{flatbuffers} but with the addition of mutable
arenas and capability references.

\section{Formal Latency Model}
\label{sec:model}

\begin{definition}[RPC latency decomposition]
Let an RPC have request size $n$ bytes spanning $f$ logical fields,
and let the network round-trip time be $\rho$ and the link bandwidth
$B$. We decompose end-to-end latency as
\[
  L \;=\; L_{\textit{ser}} \;+\; L_{\textit{net}} \;+\; L_{\textit{de}} \;+\; L_{\textit{hdl}},
\]
with transit cost $L_{\textit{net}} \;\geq\; \rho + n/B$.
\end{definition}

\subsection{Per-format costs}
Under the unit-cost RAM model with constant cache latency $c$, we have
the following bounds (justified in detail in the companion proof
document):

\begin{center}
\begin{tabular}{lll}
\toprule
Format & $L_{\textit{ser}}$ & $L_{\textit{de}}$ \\
\midrule
JSON over HTTP/2          & $\Theta(n + f \log f)$ & $\Theta(n)$ \\
Protobuf over HTTP/2      & $\Theta(f + n)$         & $\Theta(f + n)$ \\
Protobuf over QUIC        & $\Theta(f + n)$         & $\Theta(f + n)$ \\
ZAP (zero-copy)           & $\Theta(1) + \Theta(p)$ & $\Theta(1)$ \\
\bottomrule
\end{tabular}
\end{center}

\noindent
where $p$ is the number of variable-length payload bytes that must be
\emph{copied into} the arena. For fixed-layout structs $p=0$ and
$L_{\textit{ser}}=O(1)$.

\begin{remark}
The JSON $f \log f$ term arises from key sorting in canonical
JSON. The Protobuf $f$ term is varint tag decoding plus
length-prefixed dispatch. ZAP has neither because field positions are
determined at compile time.
\end{remark}

\section{Throughput and Tail-Latency Theorems}
\label{sec:theorems}

We state results here; full proofs appear in the companion proof
document at \texttt{proofs/zero-copy-throughput/proof.tex}.

\begin{theorem}[Throughput bound for zero-copy wire]
\label{thm:throughput}
Let a workload offer requests of size $n$ bytes at offered load
$\lambda$, with link bandwidth $B$ and per-request fixed CPU cost
$\sigma$ (parsing, framing, dispatch). Then the achievable throughput
$T$ for a wire format with serialization cost $L_{\textit{ser}}(n,f)$ and
deserialization cost $L_{\textit{de}}(n,f)$ satisfies
\[
  T \;\leq\;
  \min\!\Big(
    \tfrac{B}{n+h},\;
    \tfrac{1}{\sigma + L_{\textit{ser}}(n,f) + L_{\textit{de}}(n,f)}
  \Big),
\]
where $h$ is the per-message header overhead. For ZAP under fixed
layout, $L_{\textit{ser}}=L_{\textit{de}}=O(1)$ and $h$ is $O(1)$ words, so
$T \to B/n$ as $n \to \infty$. For Protobuf-over-HTTP/2,
$h \geq h_{\text{HPACK}} + h_{\text{frame}}$ and the CPU bound binds first
at small $n$, capping $T$ at $0.6$--$0.7\,B/n$.
\end{theorem}

\begin{theorem}[Tail-latency independence from field count]
\label{thm:tail}
Let $L^{(q)}$ denote the $q$-th quantile of single-RPC latency under
i.i.d.\ payloads with $f$ fields and bounded payload bytes $p$. For ZAP,
$L^{(q)}(f,p) = L^{(q)}(p) + O(1)$, i.e.\ tail latency is asymptotically
independent of field count. For Protobuf and JSON, $L^{(q)}$ grows
linearly in $f$, with the slope at the 99.9th percentile up to
$5\times$ the median slope due to allocator and HPACK
dynamic-table tail behavior~\cite{dean2013tail,xu2024microsecond}.
\end{theorem}

\section{Benchmark Methodology}
\label{sec:method}

\subsection{Hardware and software baseline}
\begin{itemize}
  \item CPU: commodity 8-core x86\_64 at 3.6\,GHz, AVX2, SMT off.
  \item Memory: 64\,GiB DDR4-3200, 2 channels, NUMA pinned.
  \item NIC: 10\,GbE, MTU 9000, RX/TX coalescing tuned for tail.
  \item OS: Linux 6.x with \texttt{tickless} and \texttt{nohz\_full}
        on benchmark cores; turbo and C-states disabled.
  \item TLS / KEM: X25519 plus a post-quantum KEM, handshake
        amortized over $10^5$ RPCs per connection.
\end{itemize}

\subsection{Workloads}
\begin{description}
  \item[Small] 16--128 byte payloads, 4--16 fields. Models cache lookups
    and feature-flag fetches.
  \item[Medium] 1--16\,KiB, 32--256 fields. Models row reads and
    structured event records.
  \item[Large] 64\,KiB--1\,MiB, lists of structs. Models bulk fetches.
  \item[Streaming] continuous server-push at 100\,k msg/s sustained for
    5 minutes; measures GC and allocator pressure.
\end{description}

\subsection{Measurement protocol}
\begin{enumerate}
  \item Warm-up: 30 seconds at target load, results discarded.
  \item Measurement window: 300 seconds.
  \item Per-RPC timestamps captured with \texttt{rdtsc} (calibrated)
        on issuer thread; clock-skew corrected per connection.
  \item Report median, $p95$, $p99$, $p999$, and max.
  \item Repeat 5 runs; report median across runs and 95\%\ CI.
  \item Congestion test: inject competing 5\,Gbps iperf flow on the
        same NIC; remeasure tails.
\end{enumerate}

\subsection{Compared protocols}
gRPC over HTTP/2 with Protobuf~\cite{grpc,rfc9113}, raw HTTP/2 with
JSON, Protobuf over QUIC~\cite{rfc9000}, and native ZAP. Each
implementation is the canonical open-source reference and is run with
its own client. Allocators are pinned (\texttt{jemalloc}) for fairness.

\section{Predicted Results}
\label{sec:results}

\Cref{tab:predicted} summarizes the headline numbers. Speedups are
ratios of the comparator's $p99$ latency to ZAP's $p99$ latency.
Numbers are calibrated against published Cap'n Proto vs Protobuf
ratios~\cite{capnproto-bench} and warehouse-scale tail
data~\cite{xu2024microsecond}.

\begin{table}[h]
\centering
\caption{Predicted $p99$ latency speedup of ZAP over comparators.}
\label{tab:predicted}
\begin{tabular}{lcccc}
\toprule
Workload & vs JSON/H2 & vs gRPC/H2 & vs PB/QUIC & ZAP $p99$ ($\mu$s) \\
\midrule
Small (64\,B, 8 fields)  & $9$--$14\times$ & $3.5$--$6\times$ & $2.5$--$4\times$ & $8\pm1$ \\
Medium (4\,KiB, 64 f.)   & $7$--$11\times$ & $4$--$7\times$   & $3$--$5\times$   & $22\pm3$ \\
Large (256\,KiB)         & $3$--$5\times$  & $1.6$--$2.2\times$ & $1.3$--$1.7\times$ & $310\pm15$ \\
Streaming 100\,k msg/s   & $6$--$9\times$  & $3$--$4\times$   & $2.5$--$3.5\times$ & --- \\
\bottomrule
\end{tabular}
\end{table}

\subsection{Tail behavior under congestion}
We expect Protobuf/HTTP-2 $p999$ to inflate by $4$--$8\times$ under the
competing iperf flow due to HPACK dynamic-table eviction and
HOL blocking. ZAP over QUIC streams should inflate by $1.5$--$2.5\times$.

\section{When ZAP Does Not Win}
\label{sec:caveats}

We are obliged to report negative regimes:
\begin{itemize}
  \item \textbf{Tiny single-field messages.} A 1-field 8-byte payload
        is dominated by transport headers; the per-message constant in
        ZAP (segment header, root pointer) costs about 16 bytes that
        Protobuf can compress with HPACK if reused.
  \item \textbf{Environments without zero-copy mmap.} On platforms that
        cannot expose contiguous arenas (some embedded RTOSes, certain
        WebAssembly runtimes pre-memory64), ZAP loses its main
        advantage and falls back to a conventional copy.
  \item \textbf{Highly heterogeneous schemas.} If consumers share the
        wire but not the schema, ZAP's positional layout is fragile
        across versions; protobuf's tag-based skip-on-unknown is
        easier to live with operationally.
  \item \textbf{Small-bandwidth high-RTT links.} When $L_{\textit{net}}$
        dominates ($\rho \gg L_{\textit{ser}}+L_{\textit{de}}$), wire format choice
        is nearly invisible.
\end{itemize}

\section{Discussion: PQ Handshake Amortization}
\label{sec:pq}

A post-quantum KEM handshake (X25519 hybridized with ML-KEM-768)
costs roughly $1$--$2$\,ms once. Amortized over a connection that
serves $10^5$ RPCs, the per-RPC overhead is $10$--$20$\,ns, well below
ZAP's per-RPC fixed cost. Long-lived ZAP connections therefore eat
the PQ handshake exactly once, in contrast to short-lived gRPC
connections that may pay it on every TLS resume failure.

\section{Related Work}
\label{sec:related}

Cap'n Proto~\cite{capnproto} originated the zero-copy arena design.
FlatBuffers~\cite{flatbuffers} adopted the same idea with a
read-only emphasis. Protocol Buffers~\cite{protobuf} and
Avro~\cite{avro} pay tag-length-value cost in exchange for
schema-evolution flexibility. MessagePack~\cite{msgpack} is JSON-like
but binary, with similar parsing cost.
HTTP/2~\cite{rfc9113} and QUIC~\cite{rfc9000} attack the transport
layer but leave the encoding layer to the application.
FaSST~\cite{kalia2016fasst} and Snap~\cite{marty2019snap}
demonstrate that, with kernel-bypass and careful encoding, RPC
latency can reach single-digit microseconds; our work targets the
encoding contribution to that budget.
The classical parser-cost analysis of Aho and Ullman~\cite{ahoullman1972}
underpins our serialization-cost lower bounds.
The ``numbers every programmer should know'' list~\cite{deannumbers}
calibrates our model's constants.

\section{Conclusion}
\label{sec:conclusion}

The serialization tax is real, measurable, and avoidable. ZAP's
zero-copy wire format collapses serialization and deserialization to
constant time in field count, which translates to predictable tail
latency and near-link-rate throughput. We have presented a formal
model, two theorems with proofs in a companion document, and a
reproducible benchmark methodology. The remaining work is empirical:
running these benchmarks across a representative hardware matrix and
publishing the raw distributions.

\printbibliography

\end{document}
